\section{\sysname System Overview}
\label{sec:overview}




The key idea of \sysname lies in leveraging a small VLM to efficiently structure video streams into discrete \emph{events}, then linking these events by extracting insights from each to construct a comprehensive index. Given a specific query, \sysname leverages this index to proactively retrieve relevant information from both the index and the associated raw frames. Ultimately, the retrieved data are utilized by the VLM to produce a coherent and contextually appropriate response.




To build \sysname as the L4 video analytics system, we established the following design principles: 1) The analytics should be scalable to any volume of video data, \ie exceeding hundreds of hours, while ensuring that the computational overhead remains independent of the video length; 2) The index construction must operate in near-real-time, allowing the system to support timely event analytics; 3) The system should accommodate not only fact-based retrieval queries but also query-focused summarization and multi-hop queries, supporting open-ended analytics.




To this end, as depicted in Fig.~\ref{fig:overview}, \sysname system is composed of two primary components: \emph{near-real-time index construction} and \emph{agentic retrieval and generation}. Within each component, we introduce a set of techniques designed to effectively realize the established design principles.





In the index construction phase, our objective is to design an effective index while ensuring construction efficiency. To achieve this, we introduce the \emph{event knowledge graph} (EKG) to structure video streams (\S\ref{subsec:event_knowledge_graph}). An EKG is a specialized form of a knowledge graph (KG) designed to represent and organize at the granularity of events and their interconnections. Recognizing that events unfold across varying temporal scales, we propose \emph{semantic chunking} (\S\ref{subsec: sematic_chunking}) to extract meaningful events from video streams. Specifically, video streams are segmented into small, uniform chunks (\eg 3-second intervals), and a small VLM, such as Qwen2.5-VL-7B, is periodically employed to generate detailed content descriptions for these chunks using carefully crafted prompts. Subsequently, neighboring chunks are merged into larger semantic chunks by identifying semantically equivalent textual descriptions with BERTScore~\cite{devlin2019bert}. For each semantic chunk, the small VLM extracts entities and their relationships. Identical entities across different events are linked to ensure consistency and coherence. Ultimately, \sysname facilitates the continuous construction of an EKG for a given video stream, regardless of its length, providing a comprehensive representation of semantic events, entities, and their interrelations in near-real-time on typical edge servers.





In the retrieval and generation process, we aim to leverage the constructed index for efficiently retrieving essential and minimal information, and to robustly generate the final answer based on the retrieved data. To achieve this, we first introduce the concept of \emph{tri-view retrieval}. Specifically, a given query undergoes simultaneous retrieval across three dimensions: events, entities, and visual embeddings. This approach ensures the acquisition of comprehensive and relevant information pertaining to the query. To further support complex queries in L4 video analytics systems \eg query-focused summarization and multi-hop queries, we propose an \emph{agentic searching} mechanism. In particular, by utilizing the LLM as an agent, \sysname proactively explores to retrieve additional information from events linked to those retrieved in earlier steps. \sysname explores multiple pathways to gather information and formulates a response to the query based on the collected data. Finally, we introduce the \emph{thoughts-consistency} strategy, which selects the most coherent and accurate final answer from multiple generated candidates.



